\documentclass[10pt,a4paper]{amsart}
\usepackage[margin=19mm]{geometry}
\usepackage{amsmath,amssymb,mathtools}
\usepackage[scr=rsfs,cal=euler]{mathalfa}
\usepackage{xfrac}
\usepackage[unicode,hidelinks,bookmarks=false]{hyperref}
\newcommand{\ie}{i.e.\ }
\newcommand{\cf}{cf.\ }
\newcommand{\resp}{respectively\ }
\newcommand{\Subsection}{\subsection}
\providecommand{\nobreakdash}{\nobreak\hbox{-}}
\setlength{\emergencystretch}{2em}
\multlinegap0pt
\usepackage{amssymb}
\usepackage{mathtools}
\newtheorem{lemma}{Lemma}
\renewcommand{\i}{\mspace{1.5mu}\mathbf i}
\renewcommand{\j}{\mspace{1.5mu}\mathbf j}
\newcommand{\n}{\mathbf n}
\renewcommand{\phi}{\varphi}
\title{The Bénard--Conway invariant of two-component links}
\author{Zedan Liu, Nikolai Saveliev}
\date{Source: arXiv:2408.16161v2 (CC BY 4.0)}
\begin{document}
\maketitle

\noindent\emph{Source and licence.} The Bénard--Conway invariant of two-component links. Version arXiv:2408.16161v2 (CC BY 4.0), distributed by arXiv under the Creative Commons Attribution 4.0 International licence, http://creativecommons.org/licenses/by/4.0/, as recorded in the arXiv licence metadata for this exact version and shown on the abstract page https://arxiv.org/abs/2408.16161v2. Full licence notice: \texttt{licenses/arxiv-2408.16161-CC-BY-4.0.txt}.\par\smallskip

\noindent\textit{Editorial scope.} Counted unit: the article's lemma L:four (source0.tex lines 596--598). The article's complete original proof of it is delivered with it (source0.tex lines 600--644, one \texttt{proof} environment of 45 lines). No mathematics is written, repaired or re-derived: the statement and the proof are the article's own lines, in the article's own notation, and no retained line is substituted or re-flowed.  Retained uncounted as the source-local context the proof needs: lines 505--507; lines 512--514; lines 558--560.  Nothing was omitted from any retained span, so the package carries no omission record.  No retained line contains a citation, so the wrapper carries no bibliography and no bibliography entry.  No retained line contains a figure environment, an image-inclusion command or drawing code, so no image is delivered and the package declares no asset dependency.  Editorial formatting only, in addition: an AMS \texttt{amsart} preamble that restates the article's own declarations and macros used by the retained lines, namely the environments \texttt{lemma} on separate counters, together with the standard packages those lines need.  The retained environments print in this document as Lemma 1 in order of appearance, whereas the article prints the same items with the numbering of its own full document.  The complete native source of the article as distributed in the arXiv e-print for this exact version is bundled unchanged as \texttt{sources/164427/source0.tex}.
\begin{equation}\label{eq:parametrize-Gamma-1}
\cos(\theta)\; =\; \frac{(|\n|^2 P_1 - d\n) \cdot Q_1}{\sin^2(\alpha_2) \sin^2(\phi)}.
\end{equation}

\begin{lemma}\label{L:one}
The right hand side of formula \eqref{eq:parametrize-Gamma-1} is a polynomial in $\cos (\phi)$, which will be denoted by $P(\cos(\phi))$.
\end{lemma}

\begin{lemma}\label{L:two}
The polynomial $P(\cos(\phi))$ has degree $2\ell$ and its leading coefficient equals $2^{2\ell -1} \sin^{2\ell} (\alpha_1) \sin^{2\ell}(\alpha_2)$.
\end{lemma}

\begin{lemma}\label{L:four}
The polynomial $P(\cos(\phi))$ evaluates to, respectively, $\cos(2 \ell(\alpha_1 + \alpha_2))$ and $\cos(2\ell(\alpha_1 - \alpha_2))$ at $\phi = 0$ and $\phi = \pi$.
\end{lemma}

\begin{proof}
We will use formula \eqref{eq:parametrize-Gamma-1}, which defines $P(\cos(\phi))$ for $\phi \in (0,\pi)$, to show that $P(\cos(\phi))$ limits to, respectively, $\cos(2\ell(\alpha_1+\alpha_2))$ and $\cos(2\ell(\alpha_1-\alpha_2))$ as $\phi \to 0$ and $\phi \to \pi$. We first calculate the limit as $\phi \to 0$. Once we eliminate the factors $\sin^2(\phi)$ as in Lemma \ref{L:one}, the calculation reduces to evaluating the expression 
\[
\frac{A_1 \cdot B_1 + A_2 \cdot B_2 + A_3 \cdot B_3}{\sin^2(\alpha_2)}
\]
at $\phi = 0$. Here, we use the notation from the proof of Lemma \ref{L:two}. It is easy to see that $A_1$, $A_2$, $A_3$ evaluate to 
\[
\begin{aligned}
A_1 & = -  \sin(\alpha_1)\sin(\alpha_2)\cos(\alpha_1+\alpha_2), \\
A_2 & =  \quad \cos(\alpha_1)\sin(\alpha_2)\sin(\alpha_1+\alpha_2), \\
A_3 & =  \quad \sin(\alpha_1)\sin(\alpha_2)\sin(\alpha_1+\alpha_2).	
\end{aligned}
\]
To evaluate $B_1, B_2$ and $B_3$, we will keep track of $\sin(\phi)$ in the formula $Q_1 = (X_1X_2)^{\ell}\cdot P_1 \cdot (X_1X_2)^{-\ell}$ while setting $\cos(\phi)$ equal to one. An induction on $\ell$ can be used to show that
\begin{equation}\label{eq:X1X2_power}
(X_1X_2)^{\ell}\, =\, q^{\ell} \,+\, \sin(\phi)\sin(\alpha_1)\,e^{\i \alpha_1}\left(q^{-\ell}\, \frac{1 - q^{2\ell}}{1 - q^2}\right)\j\, +\, \dots
\end{equation}
where $q = e^{\i(\alpha_1+\alpha_2)}$ and the dots stand for higher degree polynomials in $\sin(\phi)$. Using
the identity
\[
\frac{1 - q^{2\ell}}{1 - q^2}\, =\, e^{\i(\ell-1)(\alpha_1+\alpha_2)}\cdot\frac{\sin\left(\ell\left(\alpha_1+\alpha_2\right)\right)}{\sin(\alpha_1+\alpha_2)},
\]

\smallskip\noindent
the above formulas can be written in trigonometric form. It now follows from the formula  
\[
Q_1\, =\, (X_1X_2)^{\ell}\cdot P_1 \cdot (X_1X_2)^{-\ell}\, =\, \begin{pmatrix}
\;\qquad\quad\; B_1 \\
\;\sin(\phi) \cdot B_2 \\
\;\sin(\phi) \cdot B_3
\end{pmatrix}
\]
that, when $\phi = 0$, we have
\begin{align*}
B_1 &= 1 \\
B_2 &= \cos(2\ell(\alpha_1+\alpha_2)) + 2 \sin(\alpha_1)\big(-\sin(\alpha_2)\cos(\ell(\alpha_1+\alpha_2))\hspace{1in} \\ & \tag*{$ + \cos(\alpha_2)\sin(\ell(\alpha_1+\alpha_2))\big)\,\dfrac{\sin\left(\ell\left(\alpha_1+\alpha_2\right)\right)}{\sin(\alpha_1+\alpha_2)}$\qquad} \\
B_3 &= \sin(2\ell(\alpha_1+\alpha_2))-2\sin(\alpha_1)\big(\cos(\alpha_2)\cos(\ell(\alpha_1+\alpha_2)) \\ &
\tag*{$ + \sin(\alpha_2)\sin(\ell(\alpha_1+\alpha_2))\big)\,\dfrac{\sin\left(\ell\left(\alpha_1+\alpha_2\right)\right)}{\sin(\alpha_1+\alpha_2)}$\qquad}
\end{align*}
Finally, a tedious but straightforward trigonometric calculation using the above formulas shows that
\[
A_1 B_1 + A_2 B_2 + A_3 B_3 \, = \, \sin^2(\alpha_2)\cdot \cos(2\ell(\alpha_1+\alpha_2)),
\]
which immediately implies that $P(\cos(\phi))$ limits to $\cos(2\ell(\alpha_1+\alpha_2))$ as $\phi \to 0$. The calculation of the limit as $\phi \to \pi$ is similar.
\end{proof}

\end{document}
